\clearpage
\fancyhead[L]{\small GAUSS AN VON LINDENAU.}
\hypertarget{native-306-lower}{}
\begin{center}GAUSS AN VON LINDENAU.
\end{center}

\begin{center}
\rule{4cm}{0.4pt}\\
\fontsize{10}{12}\selectfont Monatliche Correspondenz. Band XVIII. S. 83..86. 1808 Juli.\\
\rule{4cm}{0.4pt}
\end{center}

\begin{flushright}
Göttingen, am 27. Juni 1808.
\end{flushright}

\noindent —\;—\;— Der Druck meines Werkes, obgleich es schon im Messcatalog stand, ist leider bis jetzt
noch nicht einmal zur Hälfte vollendet. Es geht damit sehr langsam, da ich mir die einzelnen
B\SourceLeaf{307}ogen zur letzten Revision zuschicken lasse, Bei Gelegenheit einer Aufgabe, die den
Gegenstand eines Abschnittes dieses Werkes ausmacht, nemlich aus \emph{vier}
geocentrischen Örtern eines Planeten, (wovon zwei unvollständig sein können) dessen
Bahn zu bestimmen, habe ich, um ein recht ausgesuchtes Beispiel zu geben, schon vor
langer Zeit noch eine Berechnung der \emph{Vesta}-Bahn gemacht, wobei eine der letzten
Mailänder Beobachtungen mit benutzt ist. Obgleich nun der Natur der Sache nach diese
Bahn nicht das möglich genaueste Resultat der vorhandenen Beobachtungen sein konnte
und sollte (weil alle angewandte Örter \emph{einzelne} Beobachtungen sind, ohne vorher
von dem wahrscheinlichsten Beobachtungsfehler befreit zu sein), so liess sich doch
voraussehen, dass diese Elemente viel genauer sein mussten, als die III$^{\text{ten}}$,
die im September schon $1'$ von den Beobachtungen abwichen. Da diese neuen Elemente
bisher noch nicht bekannt geworden sind, so theile ich Ihnen solche hier mit.


\begin{center}\fontsize{9}{11}\selectfont\setlength{\tabcolsep}{3pt}\renewcommand{\arraystretch}{1.15}\begin{tabular}{lr}
Epoche, Meridian von Paris 1807 .\;.\;.\;.\;.\;.\;.\;.\;.\;.\;.\;.\; & $168^\circ\;10'\;47''6$\\
Tägl. mittl. tropische Bewegung .\;.\;.\;.\;.\;.\;.\;.\;.\;.\;.\;.\;.\;.\; & $978''8588$\\
\shortstack[l]{Sonnenferne für die Epoche\\siderisch ruhend} .\;.\;.\;.\;.\;.\;.\;.\; & $69^\circ\;57'\;6''5$\\
Aufsteigender Knoten .\;.\;.\;.\;.\;.\;.\;.\;.\;.\;.\;.\;.\;.\;.\;.\; & $103\;11\;57''3$\\
Excentricität .\;.\;.\;.\;.\;.\;.\;.\;.\;.\;.\;.\;.\;.\;.\;.\;.\;.\;.\;.\;.\; & $0{,}0880158$\\
Logarithmus des mittl. Abstandes .\;.\;.\;.\;.\;.\;.\;.\;.\;.\;.\;.\; & $0{,}3728980$\\
Neigung der Bahn .\;.\;.\;.\;.\;.\;.\;.\;.\;.\;.\;.\;.\;.\;.\;.\;.\; & $7^\circ\;8'\;20''8$\\
\end{tabular}\end{center}


\par\indent Am 22. Juni haben wir hier die Vesta wieder aufgefunden. Seitdem haben wir zwar noch
keine heitere Nacht wieder gehabt; indess da der beobachtete Stern fast genau auf dem
Platze stand, wo die Vesta erwartet werden musste, da dieser Stern die achte Grösse
hatte (wie ungefähr der Planet jetzt sein muss), da kein anderer kenntlicher Stern in
der Nähe war, der die Vesta hätte sein können, da Hr. Professor \textsc{Harding} bei seiner
frühern Revision dieser Gegend hier keinen nicht beobachteten Fixstern bemerkt hat,
endlich da meine Beobachtungen selbst während ungefähr einer Stunde schon ziemlich
unverkennbar das Fortrücken in AR. anzuzeigen scheinen, so bleibt wohl kein Zeifel,
dass es wirklich die Vesta gewesen ist. Die erste heitere Nacht wird uns darüber volle
Gewissheit geben. Die Juno haben wir schon am 20. Juni wieder aufgefunden und
beobachtet, ob wir gleich an diesem Abend eher in einem andern Sterne 10. Grösse sie
zu erkennen glaubten, der aber am 21$^{\text{sten}}$ unverrückt seinen Platz behauptet
hatte. Am 22$^{\text{sten}}$ gab die Beobachtung volle Gewissheit. Die Juno hat kaum
die 10. Grösse, doch hoffe ich, dass Sie sie am Passagen-Instrumente beobachten werden,
da man sie 1806 bei ungefähr gleicher Lichtschwäche zu Mailand am Mauer-Quadranten
beobachtet hat.

\Needspace{10\baselineskip}
\par\indent Hier meine Beobachtungen der Juno und Vesta:


\begin{center}\fontsize{9}{11}\selectfont\setlength{\tabcolsep}{3pt}\renewcommand{\arraystretch}{1.15}\begin{tabular}{lrrrr}
& 1808 & \shortstack{Mittlere Zeit\\in Göttingen} & \shortstack{Scheinb. gerade\\Aufsteigung} & \shortstack{Scheinb. südl.\\Abweichung}\\\hline
$\JunoSymbol$ & Juni 20 & $11^{\text{u}}\,49^{\text{m}}\,0^{\text{s}}$ & $315^\circ\;29'\;34''$ & $2^\circ\;16'\;23''$\\
$\JunoSymbol$ & 22 & $12\;0\;45$ & $315\;23\;1$ & $2\;14\;20$\\
\VestaAltarThreeOhSeven & 22 & $13\;46\;25$ & $353\;12\;17$ & $9\;11\;55$\\
\end{tabular}\end{center}


\par\indent Bei der Juno geben also die 7$^{\text{ten}}$ Elemente (Gött. gelehrte Anz. 1808, St.
14, Monatl. Corresp. 1808, Febr.) die Rectascension um $11'$ zu klein, die Declination
um $2'$ zu gross; bei der Vesta geben die dritten Elemente die Rectascension um $8'$ zu
klein, die Declination um $3'$ zu gross. Die obigen Elemente hingegen geben berechnete
AR. $353^\circ\;14'\;30''$, Declination $9^\circ\;12'\;26''$, also jene um $2'\;13''$, diese
um $31''$ zu gross. Uebrigens ist die beobachtete Declination ziemlich unsicher, da der
verglichene Stern eine sehr unvortheilhafte Lage hatte.

N\SourceLeaf{308}och ein paar in diesem Monat gemachte Beobachtungen kann ich Ihnen hier mittheilen.

\par\indent Eintritt von i Virginis am dunkeln Mondsrand 4. Juni $9^{\text{u}}\,3^{\text{m}}\,39^{\text{s}}0$ M. Z.
von mir und Herrn Prof. \textsc{Harding} in Einer Secunde.

\Needspace{8\baselineskip}
\par\indent Eintritt des II$^{\text{ten}}$ Jupiter-Trabanten den 22. Juni:


\begin{center}\fontsize{10}{12}\selectfont\setlength{\tabcolsep}{3pt}\renewcommand{\arraystretch}{1.15}\begin{tabular}{ll}
$12^{\text{u}}\,45^{\text{m}}\,9^{\text{s}}$ M. Z. & \textsc{Harding} am $3\frac{1}{4}$füss. \textsc{Dollond}.\\
$12\;45\;32$\quad — & \textsc{Gauss} am 17füss. \textsc{Herschel}.\\
\end{tabular}\end{center}


\par\indent Von Hrn. Prof. \textsc{Hardings} Himmels-Charten erwarten wir nächstens das vierte Blatt.

\Needspace{16\baselineskip}
\fancyhead[L]{\small VESTA. JUNO. PALLAS. CERES.}
\begin{center}
VESTA. JUNO. PALLAS. CERES.
\end{center}

\begin{center}
\rule{4cm}{0.4pt}\\
\fontsize{10}{12}\selectfont Göttingische gelehrte Anzeigen. Stück 107. Seite 1065..1068. 1808 Juli 4.\\
\rule{4cm}{0.4pt}
\end{center}

In der Nacht vom 22. auf den 23. Juni wurde der neue Planet \emph{Vesta}, nachdem er
bei uns fast ein ganzes Jahr unsichtbar gewesen war, auf der hiesigen Sternwarte
wieder aufgefunden, und zum ersten Male wieder beobachtet. Er gleicht bereits
vollkommen einem Stern achter Grösse, und während der eine Stunde hindurch bis Tages
Anbruch fortgesetzten Beobachtungen zeigte sich schon die Ortsveränderung ganz
merklich. Das Mittel aus allen Vergleichungen mit einem Stern achter Grösse, dessen
scheinbarer Ort aus der Hist. céleste reducirt und zu


\begin{center}\fontsize{10}{12}\selectfont\setlength{\tabcolsep}{3pt}\renewcommand{\arraystretch}{1.15}\begin{tabular}{l}
$352^\circ\;44'\;1''$ gerade Aufsteigung,\\
$8^\circ\;58'\;9''$ südliche Abweichung
\end{tabular}\end{center}


\noindent angesetzt wurde, gab folgendes Resultat:


\begin{center}\fontsize{10}{12}\selectfont\setlength{\tabcolsep}{3pt}\renewcommand{\arraystretch}{1.15}\begin{tabular}{lrrr}
1808 & \shortstack{Mittlere Zeit\\in\\Göttingen} & \shortstack{Scheinb. gerade\\Aufsteigung\\der Vesta} & \shortstack{Südliche\\Abweichung\\der Vesta}\\\hline
Juni 22 & $13^{\text{u}}\,46^{\text{m}}\,25^{\text{s}}$ & $353^\circ\;12'\;17''$ & $9^\circ\;11'\;55''$\\
\end{tabular}\end{center}


\par\indent Die nach Hrn. Prof. \textsc{Gauss} dritten Elementen berechnete und im September-Heft der
Monatl. Correspondenz 1807 abgedruckte Ephemeride gibt die gerade Aufsteigung um $8'$
zu klein, die Declination um $3'$ zu gross. Jene Elemente waren schon im Juli des
vorigen Jahres berechnet, und dazu alle damals vorhandenen Beobachtungen benutzt,
wovon die letzte am 11. Juli in Lilienthal gemacht war. Die Beobachtungen der
Mailändischen Astronomen, welche den Planeten noch viel länger (bis zum 26. Sept.)
verfolgten, würden nach der Hand eine noch viel genauere Bestimmung der Bahn möglich
gemacht haben: allein obgleich diese letzten Beobachtungen nach Hrn. Prof. \textsc{Harding}'s
Berechnung sich schon über eine Minute von den dritten Elementen entfernten, so
hielt doch Herr Prof. \textsc{Gauss} damals eine neue Correction der Elemente für unnöthig,
und die schon berechnete Ephemeride zur Wiederauffindung für überflüssig genau, wie
dies nun auch der Erfolg bestätigt hat. Inzwischen hat Hr. Prof. \textsc{Gauss} nachher doch
noch eine Veranlassung zu einer neuen Berechnung der Elemente gefunden, nicht in der
Absicht, das möglich genaueste Endresultat der Beobachtungen des vorigen Jahres
auszumitteln, sondern um einen Abschnitt seines unter der Presse befindlichen Werks
über die Bahnen der Himmelskörper mit einem ausgesuchten Beispiele zu erläutern.
Hiebei wurde denn eine der letzten Mailändi\SourceLeaf{309}schen Beobachtungen mit angewandt, freilich der Absicht, warum diese Rechnung
angestellt wurde, gemäss, \emph{ohne} vorher von dem wahrscheinlichen
Beobachtungsfehler befreit zu sein. Dessen ungeachtet liess sich voraussehen, dass
das Resultat beträchtlich genauer sein würde, als die dritten Elemente: wir fügen
daher dasselbe, da es anderswo noch nicht gedruckt ist, hier bei:

\begin{center}
\emph{Elemente der Vesta.}
\end{center}


\begin{center}\fontsize{9}{11}\selectfont\setlength{\tabcolsep}{3pt}\renewcommand{\arraystretch}{1.15}\begin{tabular}{lr}
Epoche 1807 Meridian von Paris .\;.\;.\;.\;.\;.\;.\;.\; & $168^\circ\;10'\;47''6$\\
Tägliche tropische mittlere Bewegung .\;.\;.\;.\;.\;.\;.\;.\; & $978''8588$\\
Sonnennähe .\;.\;.\;.\;.\;.\;.\;.\;.\;.\;.\;.\;.\;.\;.\; & $249^\circ\;57'\;6''5$\\
Aufsteigender Knoten .\;.\;.\;.\;.\;.\;.\;.\;.\;,\;.\;.\; & $103\;11\;57$\\
\multicolumn{2}{l}{\quad beide für 1807, und siderisch ruhend}\\
Neigung der Bahn .\;.\;.\;.\;.\;.\;.\;.\;.\;.\;.\;.\;.\; & $7\;8\;21$\\
Excentricität .\;.\;.\;.\;.\;.\;.\;.\;.\;.\;,\;.\;.\;.\;.\;.\; & $0{,}0880158$\\
Logarithm des mittlern Abstandes .\;.\;.\;\hspace{.5em}.\;.\;.\;.\; & $0{,}3728980$\\
\end{tabular}\end{center}


\Needspace{9\baselineskip}
\par\indent Der nach \emph{diesen} Elementen berechnete Ort für die Zeit der obigen Beobachtung
ist:


\begin{center}\fontsize{10}{12}\selectfont\setlength{\tabcolsep}{3pt}\renewcommand{\arraystretch}{1.15}\begin{tabular}{rrrr}
gerade Aufsteigung & Unterschied & Abweichung & Unterschied\\\hline
$353^\circ\;14'\;30''$ & $+2'\;13''$ & $9^\circ\;12'\;26''$ & $+31''$\\
\end{tabular}\end{center}


\par\indent Es muss übrigens hiebei noch bemerkt werden, dass die Lage des verglichenen Sterns
zur Bestimmung der Declination sehr ungünstig war, daher diese nicht sehr
zuverlässig ist; die gerade Aufsteigung glauben wir aber um so mehr verbürgen zu
können, da von dem verglichenen Sterne zwei Beobachtungen in der Hist. cél.
vorkommen, die sehr gut übereinstimmende Resultate gaben; auf die Refraction hat
Hr. Prof. \textsc{Gauss} bei der Reduction nach einem eigenthümlichen Verfahren Rücksicht
genommen.

\par\indent Schon am 20. Juni wurde hier auch die \emph{Juno} wieder aufgefunden; sie ist noch
sehr lichtschwach, und gleicht nur einem Sterne 10..11. Grösse. Folgende
Beobachtungen haben bisher gemacht werden können: [stehen oben Seite 302 in dem
Briefe an \textsc{von Lindenau} 1808 Juni 27.]

\par\indent Die Stellung des verglichenen Sterns wurde aus der Histoire céleste bestimmt:


\begin{center}\fontsize{10}{12}\selectfont\setlength{\tabcolsep}{3pt}\renewcommand{\arraystretch}{1.15}\begin{tabular}{l}
gerade Aufsteigung $315^\circ\;24'\;58''$\\
südliche Abweichung $2^\circ\;4'\;20''$
\end{tabular}\end{center}


\par\indent Hienach weichen die siebenten, im 14. St. dieser Anzeigen mitgetheilten, Elemente
um $11'$ in gerader Aufsteigung, und um $2'$ in der Declination von der Beobachtung ab;
erstere gibt die Ephemeride (Monatl. Corresp. Februar) zu klein, letztere zu gross.
Man kann diese Differenzen schon zum Theil der Einwirkung der Störungen zuschreiben.

\par\indent Die \emph{Pallas} war schon den 20. Mai wieder aufgefunden; schlechtes Wetter,
Mondschein und zu grosse Entfernung von gut bestimmten Sternen haben aber bisher
noch keine Beobachtungen verstattet. Nach einer ungefähren Schätzung schien die
Ephemeride im December-Heft der Monatl. Corresp. die Declination gut, die
Rectascension um $2'$ zu klein zu geben.

\par\indent Die \emph{Ceres} stand bisher noch zu tief in der Morgendämmerung, wir hoffen
indess, auch von dieser bald Beobachtungen mittheilen zu können.
